\section{Troubleshooting}
\label{sec:troubleshooting}

\begin{description}
    \item[Shortcut launches but the estimation never finishes.]
        Earlier versions froze when run from the Windows shortcut
        because the \code{pythonw.exe} process inherits
        \code{sys.stdout = None}, and the worker subprocesses crashed
        in \code{tqdm}. Both the worker thread and the per-run
        subprocess now redirect a missing \code{stdout}/\code{stderr}
        to \code{os.devnull} on entry; if you still see hangs from
        the shortcut, check that you are running a recent enough
        build and report the issue with the contents of
        \code{logs/denis\_*.log}.
    \item[macOS layout looks white / clipped.]
        Old releases did not pin the Qt style and macOS rendered the
        UI with the native Aqua style, which centres tabs, paints a
        white background under the dark-themed accent colours, and
        clips widgets sized for the Fusion style. Recent releases
        force \code{app.setStyle("Fusion")} and a custom dark
        \code{QPalette} in \code{main\_window.py}; if the layout
        still looks off, raise the UI scale factor under
        \menu{File $\rightarrow$ Settings $\rightarrow$ General}.
    \item[\code{clstools} fetch fails on \code{uv sync}.]
        \code{clstools} is pulled from GitHub. If the host has no
        GitHub access, clone the repo locally and replace the entry
        in \code{[tool.uv.sources]} of \code{pyproject.toml} with a
        local path:
\begin{lstlisting}
[tool.uv.sources]
clstools = { path = "../cls_tools" }
\end{lstlisting}
        Then re-run \code{uv sync}.
    \item[Auto-Fitter never returns.]
        The Auto-Fitter spawns one worker process per CPU core
        granted in the \app{Performance} settings (default
        \code{n\_cores - 2}). If your machine is busy with another
        job, lower this number. The progress bar in the Auto-Fitter
        window also surfaces stalled workers --- if a single core
        sits at $0\,\%$ for several minutes, click \app{Stop}, drop
        the \app{N starts} value, and restart.
    \item[Model overlay disappears when switching the Pre-Analysis
        x-axis to frequency.]
        The conversion from voltage to rest-frame frequency requires
        valid laser setpoint, isotope mass, and cooler voltage. If
        any of these is missing in the source ASDF metadata, the
        axis silently falls back to MHz with no Doppler offset.
        Confirm that the \app{Cooler V} and \app{Laser cm$^{-1}$}
        fields in \app{Plot Options} are non-zero; override them
        manually if the run file is incomplete.
    \item[\code{Ctrl+C} / \code{Ctrl+V} / \code{Ctrl+Z} do nothing in
        a spinbox.]
        The application installs a global event filter for
        clipboard and undo shortcuts. If a fresh shortcut still does
        not respond, restart the application; the filter is attached
        once during startup and a partial UI build can leave it out
        of date.
    \item[Decimal separator is shown as comma instead of dot.]
        Every numeric spinbox in \denis{} is created with a
        US-English \code{QLocale}, so the decimal separator is a dot
        regardless of the operating-system locale. If a field still
        shows a comma it was created outside the shared helper ---
        report it as a bug with the field's name.
    \item[Capturing a crash.]
        Every session writes a log to
        \code{logs/denis\_<timestamp>.log} next to the application
        (a verbatim copy of the terminal output plus a structured
        log; the newest 20 files are kept). An unhandled exception
        is written under a clearly delimited \code{CRASH:} banner,
        and the startup banner records the versions needed for a bug
        report. Tick \app{Verbose session log} in
        \menu{File $\rightarrow$ Settings $\rightarrow$ General} and
        restart to raise the log level to DEBUG.
\end{description}
